\appendix

\section{Appendix}
\label{sec:appendix}

\subsection{Uploaded supporting-evidence files}

The following CSV files have been uploaded as supporting evidence for this submission and are therefore not reproduced separately in the manuscript:
\begin{itemize}
    \item \path{01_real_data_benchmark_summary.csv}
    \item \path{02_real_data_branch_ablation_summary.csv}
    \item \path{03_real_data_parameter_sensitivity_summary.csv}
    \item \path{04_real_data_tuned_candidate_summary.csv}
    \item \path{05_real_data_mechanism_compact_summary.csv}
    \item \path{06_synthetic_multiclass_benchmark_summary.csv}
\end{itemize}

\subsection{Supplementary diagnostics}

Additional fold-level rows, CKA summaries, residual traffic summaries, and full sensitivity tables are retained in the repository records directory for supplementary conversion.

\subsection{Repository identity}

The code and frozen experiment artifacts for this submission are maintained in the repository \url{https://github.com/XuelinHu/MatrixResGNNGraphClassification}. The canonical local branch used for the paper submission is \texttt{master}.

\subsection{Evidence map}

Table~\ref{tab:evidence_map} summarizes where each main empirical claim is supported in the manuscript. This compact table is included to keep the main text readable while still making the result-to-claim link explicit. In particular, the RQ3 evidence combines the class-count trend in the main text with the family-level appendix summary, so that the synthetic conclusion is supported both across class granularity and across graph-generation regimes.

\begin{table}[!htbp]
\centering
\small
\caption{Compact map from empirical claims to manuscript evidence.}
\label{tab:evidence_map}
\begin{tabular}{p{0.10\linewidth} p{0.34\linewidth} p{0.42\linewidth}}
\toprule
Research question & Claim & Primary evidence \\
\midrule
RQ1 & MatrixRes is the most frequent winner in the full benchmark but not a universal winner; GCNConv results remain dataset-dependent under a fixed operator. & Table~\ref{tab:operator_wins}, Table~\ref{tab:winners_by_operator}, Figure~\ref{fig:model_win_counts}, Table~\ref{tab:main_benchmark}, and Figure~\ref{fig:main_benchmark}. \\
RQ2 & Branch count produces a rise-then-fall pattern, and useful expansion is associated with separation without excessive residual traffic. & Table~\ref{tab:branch_best}, Figure~\ref{fig:branch_ablation}, Table~\ref{tab:mechanism_snapshot}, and Figure~\ref{fig:mechanism}. \\
RQ3 & Fold-0 sensitivity should not be treated as final evidence; matrix-style reuse becomes more favorable in higher-class synthetic structural discrimination. & Table~\ref{tab:sensitivity}, Figure~\ref{fig:sensitivity}, Table~\ref{tab:tuned_candidates}, Table~\ref{tab:synthetic_class_scaling}, Figure~\ref{fig:synthetic_multiclass_scaling}, Table~\ref{tab:synthetic_high_class}, and Table~\ref{tab:synthetic_distribution_high_class}. \\
\bottomrule
\end{tabular}
\end{table}

\subsection{Synthetic distribution-level summary}

Table~\ref{tab:synthetic_distribution_high_class} reports the higher-class synthetic results by graph family. The table is placed in the appendix because the main text focuses on the class-count scaling trend. The retained families share a common purpose: each isolates a controlled structural signal while keeping the learning pipeline unchanged. BA, ER, and the regular-degree control are degree- or density-driven settings, and MatrixRes is the strongest model in all three. SBM is community-driven and therefore imposes a different inductive pressure, but MatrixRes remains close to the best result. This pattern strengthens the central claim that local matrix residual reuse is especially useful when class labels depend on separating related structural regimes.

\begin{table}[!htbp]
\centering
\scriptsize
\caption{Synthetic high-class accuracy by graph family, averaged over \(C=4,\ldots,8\) and four operators.}
\label{tab:synthetic_distribution_high_class}
\begin{tabular}{l c c c c c}
\toprule
Family & Plain & VerticalRes & HorizontalRes & MatrixRes & MatrixResGated \\
\midrule
BA & 0.7874 & 0.8177 & 0.7855 & \(\mathbf{0.8271}\) & 0.8257 \\
ER & 0.8777 & 0.8756 & 0.8728 & \(\mathbf{0.8807}\) & 0.8796 \\
Regular & 0.7505 & 0.7923 & 0.7739 & \(\mathbf{0.8077}\) & 0.8057 \\
SBM & 0.6268 & \(\mathbf{0.6286}\) & 0.6218 & 0.6264 & 0.6176 \\
\bottomrule
\end{tabular}
\end{table}
